\element{EM4.MomentMagnetique}{
  \begin{question}{MomentMagnetique}\bareme{1}%\bareme{haut=2,p=0}
    Le moment magnétique se mesure en 
    \begin{multicols}{2}
      \begin{reponses}
        \bonne{$\SI{}{A.m^{2}}$}
        \mauvaise{$\SI{}{A.m^{-2}}$}
        \mauvaise{$\SI{}{T}$}
        \mauvaise{$\SI{}{H.m^{-1}}$}
      \end{reponses}
    \end{multicols}
  \end{question}
}
\setgroupmode{EM4.MomentMagnetique}{cyclic}


\element{EM4.Aimantation}{
  \begin{questionmult}{Aimantation}\bareme{haut=2,p=0}
    Quelles affirmations sont correctes
    \begin{multicols}{2}
      \begin{reponses}
        \bonne{L'aimantation est la densité volumique de moment magnétique}
        \mauvaise{Le moment magnétique est la densité surfacique d'aimantation}
        \bonne{Son rotationel est $\vv{j}_\text{lié}$}
        \mauvaise{Elle se mesure en $\SI{}{A.m}$}
      \end{reponses}
    \end{multicols}
  \end{questionmult}
}
\setgroupmode{EM4.Aimantation}{cyclic}


\element{EM4.Ampère}{
  \begin{question}{MaxwellAmpere}\bareme{1}
    Dans un milieu magnétique, l'équation de Maxwell-Ampère s'écrit
    \begin{multicols}{2}
      \begin{reponses}
        \bonne{$\vv{\text{rot}}\vv{H}=\vv{j}_\text{libre}$}
        \mauvaise{$\vv{\text{rot}}\vv{H}=\mu_0\vv{j}_\text{libre}$}
        \mauvaise{$\vv{\text{rot}}\vv{B}=\vv{j}_\text{libre}$}
        \mauvaise{$\vv{\text{rot}}\vv{B}=\mu_0\vv{j}_\text{libre}$}
      \end{reponses}
    \end{multicols}
  \end{question}
}
\element{EM4.Ampère}{
  \begin{question}{TheoremeAmpere}\bareme{1}
    Dans un milieu magnétique, le théorème d'Ampère s'écrit
    \begin{multicols}{2}
      \begin{reponses}
        \bonne{$\oiint\vv{H}=I_\text{libre, enlacé}$}
        \mauvaise{$\oiint\vv{H}=\mu_0I_\text{libre, enlacé}$}
        \mauvaise{$\oiint\vv{B}=I_\text{libre, enlacé}$}
        \mauvaise{$\oiint\vv{B}=\mu_0I_\text{libre, enlacé}$}
      \end{reponses}
    \end{multicols}
  \end{question}
}
\setgroupmode{EM4.Ampère}{cyclic}

\element{EM4.linearite}{
  \begin{questionmult}{linearite}\bareme{haut=2,p=0}
    Pour que $\vv{B}=\mu\vv{H}$, il est nécessaire que
    \begin{multicols}{2}
      \begin{reponses}
        \bonne{le matériau ferromagnétique soit doux}
        \mauvaise{le matériau ferromagnétique soit dur}
        \mauvaise{le milieu soit du vide}
        \mauvaise{l'on soit dans la zone de saturation}
        \bonne{l'on soit en dehors de la zone de saturation}
      \end{reponses}
    \end{multicols}
  \end{questionmult}
}
\setgroupmode{EM4.linearite}{cyclic}

\element{EM4.application}{
  \begin{questionmult}{applicationDoux}\bareme{haut=2,p=0}
    Un matériau ferromagnétique doux est utilisé pour fabriquer
    \begin{multicols}{2}
      \begin{reponses}
        \bonne{un moteur électrique}
        \bonne{un transformateur}
        \mauvaise{un aimant permanent}
      \end{reponses}
    \end{multicols}
  \end{questionmult}
}
\element{EM4.application}{
  \begin{questionmult}{applicationDur}\bareme{haut=2,p=0}
    Un matériau ferromagnétique dur est utilisé pour fabriquer
    \begin{multicols}{2}
      \begin{reponses}
        \mauvaise{un moteur électrique}
        \mauvaise{un transformateur}
        \bonne{un aimant permanent}
      \end{reponses}
    \end{multicols}
  \end{questionmult}
}
\setgroupmode{EM4.application}{cyclic}

